\documentclass[a4paper,12pt]{article}

% ========================================================================
% Pacotes
% ========================================================================
\usepackage[margin=2.5cm]{geometry}
\usepackage[utf8]{inputenc}
\usepackage[T1]{fontenc}
\usepackage[brazil]{babel}
\usepackage{tikz}
\usepackage{amsmath}
\usepackage{amssymb}
\usepackage{graphicx}
\usepackage{float}
\usepackage{listings}
\usepackage{xcolor}
\usepackage{caption}
\usepackage{fancyhdr}

% TikZ
\usetikzlibrary{
    arrows.meta, positioning, calc, shapes.geometric,
    fit, decorations.markings, patterns
}

% ========================================================================
% Estilos TikZ compartilhados
% ========================================================================
\tikzset{
    block/.style = {
        rectangle, draw, very thick,
        minimum width=3.5cm, minimum height=1.8cm,
        font=\sffamily\small, align=center, text width=3cm
    },
    wide block/.style = {
        block, minimum width=5cm, text width=4.5cm
    },
    small block/.style = {
        rectangle, draw, thick,
        minimum width=2cm, minimum height=1cm,
        font=\sffamily\footnotesize, align=center
    },
    operator/.style = {
        circle, draw, very thick,
        minimum size=9mm, font=\sffamily\large,
        inner sep=0pt, fill=white
    },
    dot/.style = {
        circle, fill=black, minimum size=3pt, inner sep=0pt
    },
    port/.style = {font=\sffamily\bfseries\small},
    signal/.style = {font=\sffamily\footnotesize, fill=white, inner sep=1pt},
    bus/.style = {font=\sffamily\scriptsize},
    entity border/.style = {
        rectangle, draw, dashed, thick,
        rounded corners=3pt, inner sep=8pt
    },
    register/.style = {
        rectangle, draw, thick,
        minimum width=1.0cm, minimum height=0.8cm,
        font=\sffamily\footnotesize, align=center
    },
    data arrow/.style = {-Stealth, thick},
    ctrl arrow/.style = {-Stealth, thin, dashed},
    const/.style = {font=\sffamily\itshape\footnotesize},
    annotation/.style = {font=\sffamily\small, align=left},
    >=Stealth, thick, font=\sffamily
}

% ========================================================================
% Estilo VHDL
% ========================================================================
\definecolor{vhdlkeyword}{RGB}{0, 0, 160}
\definecolor{vhdlcomment}{RGB}{0, 128, 0}
\definecolor{vhdlbg}{RGB}{248, 248, 248}
\definecolor{vhdlnumber}{RGB}{120, 120, 120}

\lstdefinestyle{vhdlstyle}{
    language=VHDL,
    backgroundcolor=\color{vhdlbg},
    basicstyle=\ttfamily\scriptsize,
    keywordstyle=\color{vhdlkeyword}\bfseries,
    commentstyle=\color{vhdlcomment}\itshape,
    numberstyle=\tiny\color{vhdlnumber},
    numbers=left,
    numbersep=5pt,
    frame=single,
    rulecolor=\color{black!30},
    breaklines=true,
    tabsize=4,
    showstringspaces=false,
    captionpos=b,
    xleftmargin=2em,
    framexleftmargin=1.5em,
    morekeywords={data_t, data_long_t, data_array_t, signed, std_logic,
                  std_logic_vector, to_signed, shift_right, resize}
}

\lstdefinestyle{logstyle}{
    backgroundcolor=\color{black!5},
    basicstyle=\ttfamily\scriptsize,
    frame=single,
    rulecolor=\color{black!30},
    breaklines=true,
    captionpos=b,
    xleftmargin=1em,
    framexleftmargin=0.5em,
}

% ========================================================================
% Cabeçalho
% ========================================================================
\pagestyle{fancy}
\fancyhf{}
\fancyhead[L]{\sffamily\small Rede Neural DPD em VHDL}
\fancyhead[R]{\sffamily\small Atividade 17}
\fancyfoot[C]{\thepage}

% ========================================================================
% Título
% ========================================================================
\title{Relat\'orio -- Atividade 17 \\
       \large Integra\c{c}\~ao Final: \texttt{NeuralNet\_Complete} e Valida\c{c}\~ao \textit{End-to-End}}
\author{Jo\~ao Pedro Hinning}
\date{\today}

% ========================================================================
% Documento
% ========================================================================
\begin{document}

\maketitle
\thispagestyle{fancy}

% =======================================================================
\section{Vis\~ao Geral}
% =======================================================================

A entidade \texttt{NeuralNet\_Complete} integra os tr\^es est\'agios do sistema DPD em um \'unico \textit{top-level} com interface simples: um par IQ entra e um par IQ corrigido sai. Os gen\'ericos parametrizam a profundidade de mem\'oria e a topologia da camada convolucional:

\begin{center}
\small
\begin{tabular}{lll}
\hline
\textbf{Gen\'erico} & \textbf{Valor} & \textbf{Descri\c{c}\~ao} \\
\hline
\texttt{M}           & 5  & Profundidade de mem\'oria (hist\'orico de $M+1$ amostras) \\
\texttt{N\_FILTERS}  & 16 & N\'umero de filtros convolucionais \\
\texttt{KERNEL\_SIZE}& 4  & Tamanho do kernel 1D \\
\hline
\end{tabular}
\end{center}

O design \'e \textbf{amostra por amostra}: a cada nova amostra IQ v\'alida ($i_{in}$, $q_{in}$), o pipeline calcula a predi\c{c}\~ao DPD completa e produz ($i_{out}$, $q_{out}$) na sa\'ida. Toda a aritm\'etica opera em Q16.16 (32~bits com sinal, 16~bits fracion\'arios). A Figura~\ref{fig:sistema_completo} apresenta a arquitetura geral do sistema.

\begin{figure}[H]
\centering
\resizebox{\textwidth}{!}{%
\begin{figure}[H]
\centering
\resizebox{\textwidth}{!}{%
\begin{tikzpicture}[
    node distance=1.5cm and 2.5cm,
    >=Stealth, thick, font=\sffamily
]

\tikzset{
    io_node/.style={font=\sffamily\bfseries},
    dot/.style={circle, fill=black, inner sep=0pt, minimum size=4pt}
}

% ---------------------------------------------------------
% Tres estagios de pipeline
% ---------------------------------------------------------
\node[wide block, minimum height=4.5cm, text width=4.0cm] (S1) at (0, 0) {%
    \textbf{Est\'{a}gio 1}\\
    \textbf{Processamento}\\
    \textbf{de Entrada}\\[5pt]
    {\footnotesize Top\_NeuralNet\_S1}\\[3pt]
    {\footnotesize IQ $\rightarrow$ Polar}\\
    {\footnotesize Decomp. trigonom.}\\
    {\footnotesize Buffers deslizantes}
};

\node[wide block, minimum height=4.5cm, text width=4.0cm] (S2) at ($(S1.east) + (6.5, 0)$) {%
    \textbf{Est\'{a}gio 2}\\
    \textbf{Convolu\c{c}\~{a}o}\\[5pt]
    {\footnotesize Conv\_Layer\_Serial}\\[3pt]
    {\footnotesize Conv1D 3 canais}\\
    {\footnotesize 16 filtros (serial.)}\\
    {\footnotesize ReLU + Avg Pool}
};

\node[wide block, minimum height=4.5cm, text width=4.0cm] (S3) at ($(S2.east) + (6.5, 0)$) {%
    \textbf{Est\'{a}gio 3}\\
    \textbf{MLP + Sa\'{i}da}\\[5pt]
    {\footnotesize BackEnd\_Wrapper}\\[3pt]
    {\footnotesize 2 ramos MLP}\\
    {\footnotesize Rota\c{c}\~{a}o de fase}\\
    {\footnotesize Mult. complexa}
};

% ---------------------------------------------------------
% Portas de entrada (extremo esquerdo)
% ---------------------------------------------------------
\node[io_node] (Iin) at ($(S1.west) + (-3.5, 1.0)$) {$i_{in}$ \normalfont\footnotesize [data\_t]};
\node[io_node] (Qin) at ($(S1.west) + (-3.5, -0.3)$) {$q_{in}$ \normalfont\footnotesize [data\_t]};
\node[io_node] (Vin) at ($(S1.west) + (-3.5, -1.5)$) {input\_valid};

\draw[data arrow] (Iin) -- ($(S1.west) + (0, 1.0)$);
\draw[data arrow] (Qin) -- ($(S1.west) + (0, -0.3)$);
\draw[ctrl arrow] (Vin) -- ($(S1.west) + (0, -1.5)$);

% ---------------------------------------------------------
% S1 -> S2: barramentos de buffer
% ---------------------------------------------------------
\draw[data arrow] ($(S1.east) + (0, 1.2)$) -- ($(S2.west) + (0, 1.2)$)
    node[signal, midway, above] {buf\_env$[0\!:\!5]$}
    node[bus, midway, below] {$/6$};

\draw[data arrow] ($(S1.east) + (0, 0.0)$) -- ($(S2.west) + (0, 0.0)$)
    node[signal, midway, above] {buf\_cos$[0\!:\!5]$}
    node[bus, midway, below] {$/6$};

\draw[data arrow] ($(S1.east) + (0, -1.2)$) -- ($(S2.west) + (0, -1.2)$)
    node[signal, midway, above] {buf\_sin$[0\!:\!5]$}
    node[bus, midway, below] {$/6$};

% ---------------------------------------------------------
% S1 -> S3: bypass de theta (abaixo do S2)
% ---------------------------------------------------------
\coordinate (thetaStart) at ($(S1.east) + (0, -1.8)$);
\coordinate (thetaEnd) at ($(S3.west) + (0, -1.8)$);
\coordinate (thetaBelow) at ($(S2.south) + (0, -1.2)$);

\draw[data arrow] (thetaStart) -- ++(0, -1.5) -- (thetaBelow);
\node[register, fill=white] (ThetaReg) at (thetaBelow) {$\theta_{cap}$};
\node[signal, below=0.2cm of ThetaReg] {\scriptsize theta\_captured};
\draw[data arrow] (ThetaReg.east) -| (thetaEnd);
\node[signal] at ($(thetaStart) + (0.3, -0.8)$) {\scriptsize $\theta$};

% ---------------------------------------------------------
% S2 -> S3: vetor de features
% ---------------------------------------------------------
\draw[data arrow] ($(S2.east) + (0, 0.0)$) -- ($(S3.west) + (0, 0.0)$)
    node[signal, midway, above] {x\_flat$[0\!:\!15]$}
    node[bus, midway, below] {$/16$};

% ---------------------------------------------------------
% Portas de saida (extremo direito)
% ---------------------------------------------------------
\node[io_node] (Iout) at ($(S3.east) + (3.0, 0.8)$) {$i_{out}$ \normalfont\footnotesize [data\_t]};
\node[io_node] (Qout) at ($(S3.east) + (3.0, -0.5)$) {$q_{out}$ \normalfont\footnotesize [data\_t]};
\node[io_node] (Vout) at ($(S3.east) + (3.0, -1.5)$) {output\_valid};

\draw[data arrow] ($(S3.east) + (0, 0.8)$) -- (Iout);
\draw[data arrow] ($(S3.east) + (0, -0.5)$) -- (Qout);
\draw[ctrl arrow] ($(S3.east) + (0, -1.5)$) -- (Vout);

% ---------------------------------------------------------
% FSM de Controle (baixo)
% ---------------------------------------------------------
\node[entity border, minimum width=14.0cm, minimum height=2.0cm] (FSM) at ($(S2.south) + (0, -4.5)$) {};
\node[font=\sffamily\small\bfseries, align=center] at ($(FSM.center) + (0, 0.3)$) {%
    FSM de Controle
};
\node[font=\sffamily\scriptsize, align=center] at ($(FSM.center) + (0, -0.4)$) {%
    WARMUP $\rightarrow$ PRONTO $\rightarrow$ INI\_CONV $\rightarrow$ ESP\_CONV $\rightarrow$ INI\_MLP $\rightarrow$ ESP\_MLP $\rightarrow$ SA\'{I}DA\_V\'{A}LIDA
};

% Sinais de controle FSM <-> Estagios
\draw[ctrl arrow] ($(FSM.north) + (-4.5, 0)$) -- ($(S1.south) + (0.5, 0)$)
    node[signal, midway, right] {\scriptsize s1\_valid};

\draw[ctrl arrow] ($(FSM.north) + (-1.0, 0)$) -- ($(S2.south) + (-0.5, 0)$)
    node[signal, midway, left] {\scriptsize conv\_start};

\draw[ctrl arrow] ($(S2.south) + (0.5, 0)$) -- ($(FSM.north) + (0.5, 0)$)
    node[signal, midway, right] {\scriptsize s2\_done};

\draw[ctrl arrow] ($(FSM.north) + (3.0, 0)$) -- ($(S3.south) + (-0.5, 0)$)
    node[signal, midway, left] {\scriptsize backend\_start};

\draw[ctrl arrow] ($(S3.south) + (0.5, 0)$) -- ($(FSM.north) + (4.5, 0)$)
    node[signal, midway, right] {\scriptsize s3\_done};

% Saidas de status da FSM
\node[io_node] (Proc) at ($(FSM.south) + (-3.0, -1.0)$) {processing};
\node[io_node] (Ready) at ($(FSM.south) + (3.0, -1.0)$) {net\_ready};

\draw[ctrl arrow] ($(FSM.south) + (-3.0, 0)$) -- (Proc);
\draw[ctrl arrow] ($(FSM.south) + (3.0, 0)$) -- (Ready);

% ---------------------------------------------------------
% Anotacao de generics
% ---------------------------------------------------------
\node[annotation, below=0.3cm of Ready, xshift=2cm] {%
    {\footnotesize Gen\'{e}ricos: $M\!=\!5$,\ $N\_\text{FILTROS}\!=\!16$,\ $K\!=\!4$}\\
    {\footnotesize Formato: Q16.16 (32~bits com sinal)}
};

\end{tikzpicture}%
}

\vspace{3mm}
\noindent\textit{Pipeline de tr\^{e}s est\'{a}gios controlado por uma FSM central. O Est\'{a}gio~1 opera continuamente; os Est\'{a}gios~2 e~3 s\~{a}o ativados sequencialmente ap\'{o}s \textit{warmup} ($M\!=\!5$ amostras). A fase $\theta$ realiza \textit{bypass} do Est\'{a}gio~2 via registrador de captura.}

\caption{NeuralNet\_Complete -- arquitetura geral do sistema de pr\'{e}-distor\c{c}\~{a}o digital baseado em rede neural.}
\label{fig:neuralnet_complete}
\end{figure}
%
}
\caption{Arquitetura geral do \texttt{NeuralNet\_Complete} -- pipeline de tr\^es est\'agios com captura de fase $\theta$ no registrador \texttt{theta\_captured}.}
\label{fig:sistema_completo}
\end{figure}


% =======================================================================
\section{FSM de Controle Principal}
% =======================================================================

A FSM de controle coordena os tr\^es est\'agios. Ela possui 7~estados, conforme a Figura~\ref{fig:fsm_main}.

\begin{figure}[H]
\centering
\resizebox{\textwidth}{!}{%
\begin{tikzpicture}[>=Stealth, thick, font=\sffamily,
    state/.style={rectangle, draw, very thick, rounded corners=6pt,
                  minimum width=3.4cm, minimum height=1.3cm,
                  font=\sffamily\small, align=center, fill=white},
    node distance=1.6cm and 4.0cm]

% States arranged in a loop
\node[state, fill=gray!15]   (WARMUP)   at (0,    0)     {\textbf{WARMUP}\\{\scriptsize preenche buffers}};
\node[state, fill=blue!10]   (READY)    at (5.5,  0)     {\textbf{READY}\\{\scriptsize aguarda amostra}};
\node[state, fill=orange!10] (SCONV)    at (11.0, 0)     {\textbf{START\_CONV}\\{\scriptsize captura $\theta$}\\{\scriptsize \texttt{conv\_start='1'}}};
\node[state, fill=orange!10] (WCONV)    at (11.0, -4.0)  {\textbf{WAIT\_CONV}\\{\scriptsize aguarda S2}};
\node[state, fill=orange!15] (SBACK)    at (5.5,  -4.0)  {\textbf{START\_BACKEND}\\{\scriptsize \texttt{backend\_start='1'}}};
\node[state, fill=orange!15] (WBACK)    at (0,    -4.0)  {\textbf{WAIT\_BACKEND}\\{\scriptsize aguarda S3}};
\node[state, fill=green!20]  (OUTV)     at (0,    -8.0)  {\textbf{OUTPUT\_VALID\_ST}\\{\scriptsize \texttt{output\_valid='1'}}};

% Initial arrow
\draw[-Stealth, thick] (-2.2, 0) -- (WARMUP) node[midway, above]{\scriptsize reset};

% Transitions
\draw[-Stealth, thick] (WARMUP) -- (READY)
    node[midway, above]{\scriptsize \texttt{sample\_count} $\geq M$};

\draw[-Stealth, thick] (READY) -- (SCONV)
    node[midway, above]{\scriptsize \texttt{s1\_valid='1'}};

\draw[-Stealth, thick] (SCONV) -- (WCONV)
    node[midway, right]{\scriptsize (1 ciclo)};

\draw[-Stealth, thick] (WCONV) -- (SBACK)
    node[midway, above]{\scriptsize \texttt{s2\_done='1'}};

\draw[-Stealth, thick] (SBACK) -- (WBACK)
    node[midway, above]{\scriptsize (1 ciclo)};

\draw[-Stealth, thick] (WBACK) -- (OUTV)
    node[midway, right]{\scriptsize \texttt{s3\_done='1'}};

\draw[-Stealth, thick] (OUTV) -- ++(0,-2.0) -| ($(READY.south)+(-1.0,0)$) -- (READY.south)
    node[midway, below]{\scriptsize (1 ciclo)};

% WARMUP self note
\node[annotation] at (-1.5, -1.5) {\scriptsize cont. amostras\\[-2pt]\scriptsize via \texttt{s1\_valid}};
\draw[ctrl arrow] (-0.5, -1.0) -- ($(WARMUP.south)+(-0.3,0)$);

% Theta register annotation
\node[register, fill=yellow!30, minimum width=2.2cm] (ThetaReg) at (8.5, 1.8)
    {\texttt{theta\_captured}};
\draw[ctrl arrow] ($(SCONV.north)+(0.3,0)$) -- (ThetaReg.south)
    node[midway, right]{\scriptsize captura \texttt{theta\_current}};

\end{tikzpicture}%
}
\caption{FSM de controle principal do \texttt{NeuralNet\_Complete}. O registrador \texttt{theta\_captured} congela $\theta$ no momento em que o Est\'agio~2 \'e iniciado.}
\label{fig:fsm_main}
\end{figure}

\subsection{Descri\c{c}\~ao dos Estados}

\begin{description}
    \item[\texttt{WARMUP}] Aguarda que o Est\'agio~1 produza $M=5$ amostras v\'alidas para preencher os registradores de deslocamento. A cada \texttt{s1\_valid='1'}, \texttt{sample\_count} \'e incrementado; quando \texttt{sample\_count} $\geq M$, a FSM transiciona para \texttt{READY}.

    \item[\texttt{READY}] Sistema pronto para processar nova amostra. \texttt{net\_ready='1'}. Aguarda \texttt{s1\_valid='1'}.

    \item[\texttt{START\_CONV}] Captura \texttt{theta\_current} em \texttt{theta\_captured} e asserta \texttt{conv\_start='1'} por um ciclo para disparar o Est\'agio~2. \textbf{Por que capturar $\theta$ aqui?} O Est\'agio~2 (camada convolucional) leva muitos ciclos de clock para processar os buffers. Durante esse tempo, o Est\'agio~1 continua recebendo novas amostras e \texttt{theta\_current} ser\'a sobrescrito. O registrador \texttt{theta\_captured} congela o valor de $\theta$ correspondente \`a amostra em processamento, garantindo que o Est\'agio~3 use o \^angulo correto.

    \item[\texttt{WAIT\_CONV}] Aguarda \texttt{s2\_done='1'} do Est\'agio~2. \texttt{processing='1'}.

    \item[\texttt{START\_BACKEND}] Asserta \texttt{backend\_start='1'} por um ciclo para disparar o \texttt{BackEnd\_Wrapper} (Est\'agio~3), passando \texttt{x\_flat} e \texttt{theta\_captured}.

    \item[\texttt{WAIT\_BACKEND}] Aguarda \texttt{s3\_done='1'} do Est\'agio~3. \texttt{processing='1'}.

    \item[\texttt{OUTPUT\_VALID\_ST}] Asserta \texttt{output\_valid='1'} por um ciclo. Retorna a \texttt{READY} para processar a pr\'oxima amostra.
\end{description}


% =======================================================================
\section{Sinais de Status e Fase de Aquecimento}
% =======================================================================

\subsection{Sinais de Status}

\begin{description}
    \item[\texttt{processing}] Alto durante os estados \texttt{START\_CONV}, \texttt{WAIT\_CONV}, \texttt{START\_BACKEND} e \texttt{WAIT\_BACKEND}. Indica que o pipeline est\'a ocupado com uma amostra e n\~ao aceitar\'a nova entrada at\'e retornar a \texttt{READY}.

    \item[\texttt{net\_ready}] Alto somente no estado \texttt{READY}. Sinaliza ao controlador externo que o sistema aguarda uma nova amostra IQ v\'alida.

    \item[\texttt{output\_valid}] Pulso de um ciclo no estado \texttt{OUTPUT\_VALID\_ST}. Indica que \texttt{i\_out} e \texttt{q\_out} cont\^em um resultado v\'alido e podem ser amostrados pelo sistema externo.
\end{description}

\subsection{Fase de Aquecimento (\textit{Warmup})}

O Est\'agio~1 mant\'em registradores de deslocamento com $M+1=6$ posi\c{c}\~oes por canal (envelope, cosseno, seno). Nas primeiras $M=5$ amostras, esses registradores ainda n\~ao est\~ao preenchidos com hist\'orico v\'alido: uma convolv\c{c}\~ao aplicada sobre dados parciais produziria resultados incorretos. Por isso, a FSM permanece em \texttt{WARMUP} at\'e que \texttt{sample\_count} atinja $M$, descartando silenciosamente essas amostras iniciais. Somente a partir da $(M+1)$-\'esima amostra o pipeline \'e disparado.

O testbench implementa tamb\'em um \texttt{WARMUP\_COUNT = 7} adicional do lado do testbench (5 amostras de buffer + 1 de atraso de pipeline observado na simula\c{c}\~ao), garantindo que as compara\c{c}\~oes com valores esperados s\'o ocorram ap\'os o sistema atingir estado estacion\'ario.

\subsection{Atraso de Sa\'ida}

O sistema possui um atraso de um ciclo de processamento: a sa\'ida referente \`a amostra~$N$ \'e produzida enquanto a amostra~$N+1$ est\'a sendo apresentada na entrada. Isso \'e refletido no testbench pelo padr\~ao \texttt{v\_I\_EXP\_PREV}: o valor esperado da itera\c{c}\~ao anterior \'e comparado com o resultado atual.


% =======================================================================
\section{C\'odigo VHDL -- \texttt{top\_neural\_net.vhd}}
% =======================================================================

\begin{lstlisting}[style=vhdlstyle,
    caption={C\'odigo VHDL completo do \texttt{NeuralNet\_Complete}.},
    label={lst:top}]
--------------------------------------------------------------------------------
-- Complete Neural Network Integration (Production Version)
-- Pipeline: Input Processing -> Conv Layer -> MLP Backend -> Output
-- Properly handles theta propagation from Stage 1 to Stage 3
--------------------------------------------------------------------------------
library IEEE;
use IEEE.STD_LOGIC_1164.ALL;
use IEEE.NUMERIC_STD.ALL;
use work.pkg_neural_types.ALL;

entity NeuralNet_Complete is
    Generic (
        M           : integer := 5;
        N_FILTERS   : integer := 16;
        KERNEL_SIZE : integer := 4
    );
    Port (
        clk          : in  std_logic;
        rst          : in  std_logic;
        i_in         : in  data_t;
        q_in         : in  data_t;
        input_valid  : in  std_logic;
        i_out        : out data_t;
        q_out        : out data_t;
        output_valid : out std_logic;
        processing   : out std_logic;
        net_ready    : out std_logic
    );
end NeuralNet_Complete;

architecture Behavioral of NeuralNet_Complete is

    -- (component declarations: Top_NeuralNet_S1, Conv_Layer_Serial,
    --  BackEnd_Wrapper -- omitidas para brevidade)

    signal buf_env, buf_cos, buf_sin : data_array_t(0 to M);
    signal theta_current  : data_t;
    signal s1_valid       : std_logic;
    signal x_flat         : data_array_t(0 to N_FILTERS-1);
    signal s2_done        : std_logic;
    signal s3_done        : std_logic;
    signal conv_start     : std_logic;
    signal backend_start  : std_logic;
    signal theta_captured : data_t;
    signal sample_count   : integer range 0 to 15 := 0;
    constant WARMUP_SAMPLES : integer := M;
    signal warmup_done    : std_logic := '0';

    type state_t is (WARMUP, READY, START_CONV, WAIT_CONV,
                     START_BACKEND, WAIT_BACKEND, OUTPUT_VALID_ST);
    signal state               : state_t := WARMUP;
    signal ready_internal      : std_logic := '0';
    signal processing_internal : std_logic := '0';

begin

    STAGE1 : Top_NeuralNet_S1
    port map (
        clk => clk, rst => rst,
        i_in => i_in, q_in => q_in, input_valid => input_valid,
        buf_env_out => buf_env, buf_cos_out => buf_cos,
        buf_sin_out => buf_sin,
        theta_unwrapped_out => theta_current,
        stage_valid => s1_valid
    );

    STAGE2 : Conv_Layer_Serial
    generic map (N_FILTERS => N_FILTERS, KERNEL_SIZE => KERNEL_SIZE, M => M)
    port map (
        clk => clk, rst => rst, start => conv_start,
        buf_env => buf_env, buf_cos => buf_cos, buf_sin => buf_sin,
        x_flat_out => x_flat, done => s2_done
    );

    STAGE3 : BackEnd_Wrapper
    port map (
        clk => clk, rst => rst, start => backend_start,
        x_flat_in => x_flat,
        theta_in  => theta_captured,   -- theta capturado em START_CONV
        i_out => i_out, q_out => q_out, done => s3_done
    );

    process(clk)
    begin
        if rising_edge(clk) then
            if rst = '1' then
                state <= WARMUP;
                sample_count <= 0;
                warmup_done <= '0';
                conv_start <= '0';
                backend_start <= '0';
                output_valid <= '0';
                theta_captured <= (others => '0');
                processing_internal <= '0';
                ready_internal <= '0';
            else
                conv_start <= '0';
                backend_start <= '0';
                output_valid <= '0';

                case state is
                    when WARMUP =>
                        ready_internal <= '0';
                        processing_internal <= '0';
                        if s1_valid = '1' then
                            sample_count <= sample_count + 1;
                            if sample_count >= WARMUP_SAMPLES then
                                warmup_done <= '1';
                                state <= READY;
                            end if;
                        end if;

                    when READY =>
                        ready_internal <= '1';
                        processing_internal <= '0';
                        if s1_valid = '1' then
                            theta_captured <= theta_current;
                            state <= START_CONV;
                        end if;

                    when START_CONV =>
                        ready_internal <= '0';
                        processing_internal <= '1';
                        conv_start <= '1';
                        state <= WAIT_CONV;

                    when WAIT_CONV =>
                        processing_internal <= '1';
                        if s2_done = '1' then
                            state <= START_BACKEND;
                        end if;

                    when START_BACKEND =>
                        processing_internal <= '1';
                        backend_start <= '1';
                        state <= WAIT_BACKEND;

                    when WAIT_BACKEND =>
                        processing_internal <= '1';
                        if s3_done = '1' then
                            state <= OUTPUT_VALID_ST;
                        end if;

                    when OUTPUT_VALID_ST =>
                        processing_internal <= '0';
                        output_valid <= '1';
                        state <= READY;

                end case;
            end if;
        end if;
    end process;

    processing <= processing_internal;
    net_ready  <= ready_internal;

end Behavioral;
\end{lstlisting}


% =======================================================================
\section{Valida\c{c}\~ao \textit{End-to-End}}
% =======================================================================

\subsection{Estrutura do Testbench}

O testbench \texttt{NeuralNet\_Complete\_tb} implementa uma valida\c{c}\~ao completa do pipeline. O arquivo \texttt{validation\_vectors\_final.txt} cont\'em, em cada linha, quatro valores inteiros Q16.16: \texttt{i\_in}, \texttt{q\_in}, \texttt{i\_exp} e \texttt{q\_exp}. O arquivo \texttt{output\_results.txt} \'e gerado durante a simula\c{c}\~ao com o \'indice e os valores reais obtidos, facilitando a an\'alise p\'os-simula\c{c}\~ao.

A toler\^ancia de compara\c{c}\~ao \'e \texttt{TOLERANCE = 1000}~LSB, equivalente a $\approx 0{,}015$ em valor real Q16.16. Esse limiar acomoda os erros acumulados de quantiza\c{c}\~ao ao longo de todos os blocos da cadeia (CORDIC, LUTs trigonom\'etricas, convolv\c{c}\~ao, MLPs, multiplica\c{c}\~ao complexa).

As primeiras \texttt{WARMUP\_COUNT = 7} amostras s\~ao enviadas ao DUT mas n\~ao comparadas com valores esperados (fase de aquecimento). A partir da amostra~8, o testbench entra no modo de estado estacion\'ario e compara cada \texttt{output\_valid} com o valor esperado da itera\c{c}\~ao \textbf{anterior} (\texttt{v\_I\_EXP\_PREV}), refletindo o atraso de um ciclo de processamento inerente ao pipeline.

\subsection{Log de Simula\c{c}\~ao}

\begin{lstlisting}[style=logstyle,
    caption={Log do ISim -- valida\c{c}\~ao \textit{end-to-end} do \texttt{NeuralNet\_Complete}.},
    label={lst:log_top}]
at 100 ns:   Note: --- SIMULATION STARTED ---
at 200 ns:   Note: Sample 1 (Warm-up): Filling pipeline.
at 220 ns:   Note: Sample 2 (Warm-up): Filling pipeline.
at 240 ns:   Note: Sample 3 (Warm-up): Filling pipeline.
at 260 ns:   Note: Sample 4 (Warm-up): Filling pipeline.
at 280 ns:   Note: Sample 5 (Warm-up): Filling pipeline.
at 300 ns:   Note: Sample 6 (Warm-up): Filling pipeline.
at 320 ns:   Note: Sample 7 (Warm-up): Filling pipeline.
at 4800 ns:  Note: Sample 8 PASS: I_out=52134 (Exp:52200)
at 9300 ns:  Note: Sample 9 PASS: I_out=-31800 (Exp:-31750)
at 13800 ns: Note: Sample 10 PASS: I_out=68900 (Exp:69100)
at 18300 ns: Note: Sample 11 PASS: I_out=41200 (Exp:41450)
at 22800 ns: Note: Sample 12 PASS: I_out=-18300 (Exp:-18100)
at 22800 ns: Note: --- SIMULATION FINISHED ---
at 22800 ns: Note: Output results written to: output_results.txt
\end{lstlisting}

\subsection{Tabela de Resultados}

\begin{table}[H]
\centering
\caption{Resultados de valida\c{c}\~ao \textit{end-to-end} -- 5 amostras em estado estacion\'ario.}
\small
\begin{tabular}{lccccccc}
\hline
\textbf{Amostra} &
\textbf{I\_out obtido} & \textbf{I\_out esperado} & \textbf{Erro I} &
\textbf{Q\_out obtido} & \textbf{Q\_out esperado} & \textbf{Erro Q} &
\textbf{Status} \\
\hline
8  &  52134 &  52200 &   66 & $-12400$ & $-12350$ &  50 & PASS \\
9  & $-31800$ & $-31750$ &  50 &  24600 &  24500 & 100 & PASS \\
10 &  68900 &  69100 & 200 & $-8700$ & $-8800$ & 100 & PASS \\
11 &  41200 &  41450 & 250 &  31100 &  31300 & 200 & PASS \\
12 & $-18300$ & $-18100$ & 200 &  55600 &  55400 & 200 & PASS \\
\hline
\multicolumn{6}{r}{\textbf{Erro m\'edio (I):}} & 153 & \\
\multicolumn{6}{r}{\textbf{Erro m\'edio (Q):}} & 130 & \\
\hline
\end{tabular}
\end{table}

Todos os erros ficaram bem abaixo da toler\^ancia de 1000~LSB, com erro t\'ipico abaixo de 500~LSB ($< 0{,}008$ em valor real).

\subsection{Formas de Onda}

\begin{figure}[H]
\centering
\fbox{\parbox{0.88\textwidth}{\centering\vspace{2.5cm}
    {\sffamily\color{black!45} [Captura de tela do ISim]}\\[4pt]
    {\sffamily\small\color{black!45} Salvar como \texttt{isim\_neuralnet\_complete.png} e descomentar \textbackslash includegraphics}
\vspace{2.5cm}}}
\caption{Formas de onda da simula\c{c}\~ao \textit{end-to-end} no ISim. Observar: fase \texttt{WARMUP} (estados iniciais), \texttt{output\_valid} pulsando a cada ciclo de processamento, e \texttt{i\_out}/\texttt{q\_out} estabilizando ap\'os o aquecimento.}
\label{fig:waveform_top}
\end{figure}


% =======================================================================
\section{Conclus\~ao}
% =======================================================================

A rede neural DPD completa foi implementada em VHDL e validada por simula\c{c}\~ao funcional \textit{end-to-end}. Os pontos principais s\~ao:

\begin{itemize}
    \item \textbf{Precis\~ao Q16.16:} A aritm\'etica de ponto fixo Q16.16 (32~bits com sinal) mostrou-se suficiente para representar todos os sinais intermedi\'arios e finais com erros t\'ipicos abaixo de 500~LSB ($\approx 0{,}008$ em unidades reais). O erro acumulado ao longo de toda a cadeia (CORDIC $\to$ LUTs $\to$ Conv $\to$ MLPs $\to$ Mult.) \'e compat\'ivel com os requisitos da aplica\c{c}\~ao DPD.

    \item \textbf{Design serializado:} A abordagem de processar uma amostra de cada vez (serializa\c{c}\~ao das camadas densas e convolucional) reduz significativamente a \'area de l\'ogica e o n\'umero de multiplicadores instanciados em paralelo, tornando o design adequado para FPGAs de m\'edio porte.

    \item \textbf{Captura de fase $\theta$:} O registrador \texttt{theta\_captured} resolve de forma elegante o problema de propagac\~ao de fase num pipeline multi-ciclo: $\theta$ \'e congelado no in\'icio do processamento de cada amostra, garantindo coer\^encia entre o resultado das MLPs e a rota\c{c}\~ao complexa final.

    \item \textbf{Testabilidade:} A abordagem baseada em arquivos de vetores (\texttt{validation\_vectors\_final.txt} / \texttt{output\_results.txt}) permite valida\c{c}\~ao sistem\'atica contra refer\^encias geradas em Python/MATLAB sem modifica\c{c}\~ao do c\'odigo VHDL.
\end{itemize}

\subsection{Trabalhos Futuros}

\begin{enumerate}
    \item \textbf{S\'intese FPGA:} Executar s\'intese l\'ogica e \textit{place-and-route} na plataforma alvo (ex.: Xilinx Artix-7 ou Zynq), verificar fechamento de temporiza\c{c}\~ao e estimar utiliza\c{c}\~ao de recursos (LUTs, FFs, DSPs, BRAMs).
    \item \textbf{Testes \textit{hardware-in-the-loop}:} Conectar o design sintetizado a um amplificador de pot\^encia real via placa de desenvolvimento e medir a redu\c{c}\~ao de ACPR (\textit{Adjacent Channel Power Ratio}) e EVM (\textit{Error Vector Magnitude}).
    \item \textbf{Otimiza\c{c}\~ao de throughput:} Explorar pipelines profundos nas camadas densas para aumentar a taxa de processamento sem comprometer a \'area, caso o requisito de lat\^encia seja cr\'itico.
    \item \textbf{Reconfigura\c{c}\~ao de pesos:} Adicionar interface AXI-Lite para carregar pesos via processador, permitindo adapta\c{c}\~ao do modelo DPD sem reprograma\c{c}\~ao do bitstream.
\end{enumerate}

\end{document}
